\SGANextSectionTarget{sga3:xxiii:section:4}\SGASection{4}{Uniqueness of Pinned Groups: Fundamental Theorem}

\SGARefTarget{sga3:XXIII:4.1}
\SGAInlineRefTarget{sga3:xxiii:theorem:4-1}\SGAStatement{Theorem 4.1}{Let \(S\) be a nonempty scheme. The
functor \(R\) of \SGARefLink{sga3:XXIII:1.6}{1.6} is fully faithful. More precisely, let \(G,G'\)
be pinned \(S\)-groups, let \(p>0\) be an integer such that
\(x\mapsto x^p\) is an endomorphism of \(\mathbf G_{a,S}\), and let
\[
  h:R(G')\longrightarrow R(G)
\]
be a \(p\)-morphism of pinned root data. There is a unique morphism
of pinned groups \(f:G\to G'\) such that \(R(f)=h\).}

Uniqueness was proved in \SGARefLink{sga3:XXIII:1.9}{1.9}, so only existence remains. By
hypothesis there are a bijection \(d:R\xrightarrow{\sim}R'\) and a
map
\[
  q:R\longrightarrow\{p^n\mid n\in\mathbf N\}
\]
such that, for every \(\alpha\in R\),
\[
  h(d(\alpha))=q(\alpha)\alpha,
  \qquad
  {}^th(\alpha^*)=q(\alpha)d(\alpha)^*.
\]
In particular, \(G\) and \(G'\) have the same semisimple rank.

\SGARefTarget{sga3:XXIII:4.1.1}
\SGAParagraph{4.1.1}
Suppose that both semisimple ranks are \(0\). Then \(G,G'\) are
tori:
\[
  G=T=D_S(M),\qquad G'=T'=D_S(M'),
\]
and \(h:M'\to M\) is an ordinary homomorphism. Take \(f=D_S(h)\).

\SGARefTarget{sga3:XXIII:4.1.2}
\SGAParagraph{4.1.2}
Suppose that both semisimple ranks are \(1\), and put
\[
  f_T=D_S(h):T\longrightarrow T'.
\]
Writing \(\alpha'=d(\alpha)\), the diagram
\[
\begin{tikzcd}[column sep=large,row sep=large]
\mathbf G_{m,S}\ar[r,"\alpha^*"]\ar[d,"q(\alpha)"]
  &T\ar[r,"\alpha"]\ar[d,"f_T"']
  &\mathbf G_{m,S}\ar[d,"q(\alpha)"]\\
\mathbf G_{m,S}\ar[r,"\alpha'^*"]
  &T'\ar[r,"\alpha'"]
  &\mathbf G_{m,S}
\end{tikzcd}
\]
commutes. Apply \SGARefLink{sga3:xx:4.1}{Exposé XX, 4.1}.

\SGARefTarget{sga3:XXIII:4.1.3}
\SGAParagraph{4.1.3}
Suppose that both semisimple ranks are \(2\). \SGARefLink{sga3:XXI:7.5.3}{Exposé XXI, 7.5.3}
lists every possible \(h\); we verify the conditions of \SGARefLink{sga3:XXIII:2.5}{Corollary 2.5}
in each case. Write
\[
  \Delta=\{\alpha,\beta\},\qquad
  \Delta'=\{\alpha',\beta'\},\qquad
  d(\alpha)=\alpha',\quad d(\beta)=\beta'.
\]

\SGARefTarget{sga3:XXIII:4.1.4}
\SGAParagraph{4.1.4. \(G\) and \(G'\) of type
\(A_1\times A_1\)}
Here
\[
  h(\alpha')=q\alpha,\qquad h(\beta')=q_1\beta.
\]
Conditions \SGARefLink{sga3:XXIII:2.5:item:ii}{2.5(ii)} and (iii) follow from \SGARefLink{sga3:xxiii:proposition:3-1-2:item:ii}{Proposition 3.1.2(ii)},(iii).
For (i),
\[
\begin{aligned}
 D_S(h)t_{\alpha\beta}
  &=D_S(h)(t_\alpha t_\beta)\\
  &={}^th(\alpha^*)(-1)\,{}^th(\beta^*)(-1)\\
  &=\alpha'^*((-1)^q)\,\beta'^*((-1)^{q_1}).
\end{aligned}
\]
The assumption on \(x\mapsto x^p\), with \(S\ne\varnothing\), implies
that \(p=1\) or \(S\) has characteristic \(p\). Thus
\((-1)^q=-1\): if \(q\) is even, then \(p=2\) and \(1=-1\).
Consequently
\[
  D_S(h)t_{\alpha\beta}
   =\alpha'^*(-1)\beta'^*(-1)=t'_{\alpha'\beta'}.
\]

\SGARefTarget{sga3:XXIII:4.1.5}
\SGAParagraph{4.1.5. \(G\) and \(G'\) of type \(A_2\)}
Here
\[
  h(\alpha')=q\alpha,\qquad h(\beta')=q\beta.
\]
Put
\[
  X_{\alpha+\beta}=\operatorname{Ad}(w_\beta)X_\alpha,
  \qquad
  X'_{\alpha'+\beta'}=\operatorname{Ad}(w'_{\beta'})X'_{\alpha'}.
\]
Condition (i) follows as above from \SGARefLink{sga3:xxiii:proposition:3-2-1:item:i}{Proposition 3.2.1(i)}, and (ii)
follows immediately from \SGARefLink{sga3:xxiii:proposition:3-2-1:item:ii}{Proposition 3.2.1(ii)}. For (iii), the map
\[
\begin{aligned}
 p_\alpha(x)p_\beta(y)p_{\alpha+\beta}(z)
 \longmapsto
 p'_{\alpha'}(x^q)p'_{\beta'}(y^q)p'_{\alpha'+\beta'}(z^q)
\end{aligned}
\]
must be a homomorphism. The only nontrivial relation is
\[
 p_\beta(y)p_\alpha(x)
  =p_\alpha(x)p_\beta(y)p_{\alpha+\beta}(xy),
\]
which maps to
\[
 p'_{\beta'}(y^q)p'_{\alpha'}(x^q)
  =p'_{\alpha'}(x^q)p'_{\beta'}(y^q)
    p'_{\alpha'+\beta'}(x^qy^q).
\]

\SGARefTarget{sga3:XXIII:4.1.6}
\SGAParagraph{4.1.6}
The same argument treats types \(B_2\) and \(G_2\) when the
radicial exponents are equal, using \SGARefLink{sga3:xxiii:proposition:3-3-1}{Propositions 3.3.1} and \SGARefLink{sga3:xxiii:proposition:3-4-1}{3.4.1}.
It remains to treat the two exceptional cases in \SGARefLink{sga3:XXI:7.5.3}{Exposé XXI, 7.5.3}.

\SGARefTarget{sga3:XXIII:4.1.7}
\SGAParagraph{4.1.7. Type \(B_2\), characteristic \(2\),
\(q(\alpha)=2q\), \(q(\beta)=q\)}
The positive roots are
\[
 \{\alpha,\beta,\alpha+\beta,2\alpha+\beta\}
\]
and
\[
 \{\alpha',\beta',\alpha'+\beta',\alpha'+2\beta'\},
\]
where \(\alpha\) and \(\beta'\) are the short simple roots. We have
\[
\begin{aligned}
 h(\alpha')&=2q\alpha,&h(\beta')&=q\beta,\\
 h(\alpha'+\beta')&=q(2\alpha+\beta),&
 h(\alpha'+2\beta')&=2q(\alpha+\beta),
\end{aligned}
\]
so
\[
\begin{aligned}
 d(\alpha+\beta)&=\alpha'+2\beta',&q(\alpha+\beta)&=2q,\\
 d(2\alpha+\beta)&=\alpha'+\beta',&q(2\alpha+\beta)&=q.
\end{aligned}
\]
Put
\[
\begin{aligned}
 X_{\alpha+\beta}&=\operatorname{Ad}(w_\beta)X_\alpha,&
 X_{2\alpha+\beta}&=\operatorname{Ad}(w_\alpha)X_\beta,\\
 X'_{\alpha'+\beta'}&=\operatorname{Ad}(w'_{\alpha'})X'_{\beta'},&
 X'_{\alpha'+2\beta'}&=\operatorname{Ad}(w'_{\beta'})X'_{\alpha'}.
\end{aligned}
\]

For (i), characteristic \(2\) gives \(-1=1\), hence
\[
 t_{\alpha\beta}=t_\alpha=\alpha^*(-1)=e
   =\beta'^*(-1)=t'_{\beta'}=t'_{\alpha'\beta'}.
\]
For (ii), the two defining identities are
\[
\begin{aligned}
 \operatorname{Ad}(w_\alpha)X_\beta
  &=X_{2\alpha+\beta},&
 \operatorname{Ad}(w'_{\alpha'})X'_{\beta'}
  &=X'_{\alpha'+\beta'}=X'_{d(2\alpha+\beta)},\\
 \operatorname{Ad}(w_\beta)X_\alpha
  &=X_{\alpha+\beta},&
 \operatorname{Ad}(w'_{\beta'})X'_{\alpha'}
  &=X'_{\alpha'+2\beta'}=X'_{d(\alpha+\beta)}.
\end{aligned}
\]
\SGARefLink{sga3:xxiii:proposition:3-3-1:item:ii}{Proposition 3.3.1(ii)}, together with \(-1=1\), gives the four
remaining identities:
\[
\begin{aligned}
 \operatorname{Ad}(w_\alpha)X_{\alpha+\beta}&=X_{\alpha+\beta},&
 \operatorname{Ad}(w'_{\alpha'})X'_{d(\alpha+\beta)}
   &=X'_{d(\alpha+\beta)},\\
 \operatorname{Ad}(w_\alpha)X_{2\alpha+\beta}&=X_\beta,&
 \operatorname{Ad}(w'_{\alpha'})X'_{d(2\alpha+\beta)}
   &=X'_{d(\beta)},\\
 \operatorname{Ad}(w_\beta)X_{\alpha+\beta}&=X_\alpha,&
 \operatorname{Ad}(w'_{\beta'})X'_{d(\alpha+\beta)}
   &=X'_{d(\alpha)},\\
 \operatorname{Ad}(w_\beta)X_{2\alpha+\beta}&=X_{2\alpha+\beta},&
 \operatorname{Ad}(w'_{\beta'})X'_{d(2\alpha+\beta)}
   &=X'_{d(2\alpha+\beta)}.
\end{aligned}
\]

For (iii), the sole nontrivial relations are
\[
\begin{aligned}
 p_\beta(y)p_\alpha(x)
  &=p_\alpha(x)p_\beta(y)p_{\alpha+\beta}(xy)
    p_{2\alpha+\beta}(x^2y),\\
 p'_{\alpha'}(y')p'_{\beta'}(x')
  &=p'_{\beta'}(x')p'_{\alpha'}(y')
    p'_{\alpha'+\beta'}(x'y')
    p'_{\alpha'+2\beta'}(x'y'^2).
\end{aligned}
\]
The proposed map is
\[
\begin{aligned}
 &p_\alpha(x)p_\beta(y)p_{\alpha+\beta}(z)p_{2\alpha+\beta}(t)\\
 &\quad\longmapsto
 p'_{\alpha'}(x^{2q})p'_{\beta'}(y^q)
 p'_{\alpha'+2\beta'}(z^{2q})p'_{\alpha'+\beta'}(t^q).
\end{aligned}
\]
It respects the first relation because the second, with
\(y'=x^{2q}\) and \(x'=y^q\), reads
\[
\begin{aligned}
 p'_{\beta'}(y^q)p'_{\alpha'}(x^{2q})
  ={}&p'_{\alpha'}(x^{2q})p'_{\beta'}(y^q)
   p'_{\alpha'+2\beta'}((xy)^{2q})\\
   &{}\cdot p'_{\alpha'+\beta'}((x^2y)^q).
\end{aligned}
\]

\SGARefTarget{sga3:XXIII:4.1.8}
\SGAParagraph{4.1.8. Type \(G_2\), characteristic \(3\),
\(q(\alpha)=3q\), \(q(\beta)=q\)}
The positive roots are
\[
\{\alpha,\beta,\alpha+\beta,2\alpha+\beta,3\alpha+\beta,
  3\alpha+2\beta\}
\]
and
\[
\{\alpha',\beta',\alpha'+\beta',\alpha'+2\beta',
  \alpha'+3\beta',2\alpha'+3\beta'\},
\]
where again \(\alpha\) and \(\beta'\) are short. We have
\[
\begin{aligned}
 h(\alpha')&=3q\alpha,&h(\beta')&=q\beta,\\
 h(\alpha'+\beta')&=q(3\alpha+\beta),&
 h(\alpha'+2\beta')&=q(3\alpha+2\beta),\\
 h(\alpha'+3\beta')&=3q(\alpha+\beta),&
 h(2\alpha'+3\beta')&=3q(2\alpha+\beta),
\end{aligned}
\]
and hence
\[
\begin{aligned}
 d(\alpha+\beta)&=\alpha'+3\beta',&q(\alpha+\beta)&=3q,\\
 d(2\alpha+\beta)&=2\alpha'+3\beta',&q(2\alpha+\beta)&=3q,\\
 d(3\alpha+\beta)&=\alpha'+\beta',&q(3\alpha+\beta)&=q,\\
 d(3\alpha+2\beta)&=\alpha'+2\beta',&q(3\alpha+2\beta)&=q.
\end{aligned}
\]
Define
\[
\begin{aligned}
 X_{\alpha+\beta}&=\operatorname{Ad}(w_\beta)X_\alpha,&
 X_{2\alpha+\beta}&=\operatorname{Ad}(w_\alpha)X_{\alpha+\beta},\\
 X_{3\alpha+\beta}&=-\operatorname{Ad}(w_\alpha)X_\beta,&
 X_{3\alpha+2\beta}&=\operatorname{Ad}(w_\beta)X_{3\alpha+\beta},\\
 X'_{\alpha'+\beta'}&=-\operatorname{Ad}(w'_{\alpha'})X'_{\beta'},&
 X'_{\alpha'+2\beta'}&=\operatorname{Ad}(w'_{\beta'})X'_{\alpha'+\beta'},\\
 X'_{\alpha'+3\beta'}&=\operatorname{Ad}(w'_{\beta'})X'_{\alpha'},&
 X'_{2\alpha'+3\beta'}&=\operatorname{Ad}(w'_{\alpha'})X'_{\alpha'+3\beta'}.
\end{aligned}
\]
Condition (i) follows from \SGARefLink{sga3:xxiii:proposition:3-4-1:item:i}{Proposition 3.4.1(i)}. Four instances of
(ii) hold by definition:
\[
\begin{aligned}
 \operatorname{Ad}(w_\beta)X_\alpha&=X_{\alpha+\beta},&
 \operatorname{Ad}(w'_{\beta'})X'_{\alpha'}&=X'_{d(\alpha+\beta)},\\
 \operatorname{Ad}(w_\alpha)X_\beta&=-X_{3\alpha+\beta},&
 \operatorname{Ad}(w'_{\alpha'})X'_{\beta'}&=-X'_{d(3\alpha+\beta)},\\
 \operatorname{Ad}(w_\beta)X_{3\alpha+\beta}&=X_{3\alpha+2\beta},&
 \operatorname{Ad}(w'_{\beta'})X'_{d(3\alpha+\beta)}
   &=X'_{d(3\alpha+2\beta)},\\
 \operatorname{Ad}(w_\alpha)X_{\alpha+\beta}&=X_{2\alpha+\beta},&
 \operatorname{Ad}(w'_{\alpha'})X'_{d(\alpha+\beta)}
   &=X'_{d(2\alpha+\beta)}.
\end{aligned}
\]
The remaining instances follow from \SGARefLink{sga3:xxiii:proposition:3-4-1:item:ii}{Proposition 3.4.1(ii)}; for
example,
\[
\begin{aligned}
 \operatorname{Ad}(w_\alpha)X_{2\alpha+\beta}&=-X_{\alpha+\beta},&
 \operatorname{Ad}(w'_{\alpha'})X'_{d(2\alpha+\beta)}
   &=-X'_{d(\alpha+\beta)},\\
 \operatorname{Ad}(w_\beta)X_{3\alpha+2\beta}&=-X_{3\alpha+\beta},&
 \operatorname{Ad}(w'_{\beta'})X'_{d(3\alpha+2\beta)}
   &=-X'_{d(3\alpha+\beta)}.
\end{aligned}
\]

For (iii), in characteristic \(3\), the nontrivial relations in
\(U\) are
\[
\begin{aligned}
 p_\beta(y)p_\alpha(x)
  ={}&p_\alpha(x)p_\beta(y)p_{\alpha+\beta}(xy)
  p_{2\alpha+\beta}(x^2y)p_{3\alpha+\beta}(x^3y)
  p_{3\alpha+2\beta}(x^3y^2),\\
 p_{\alpha+\beta}(y)p_\alpha(x)
  ={}&p_\alpha(x)p_{\alpha+\beta}(y)p_{2\alpha+\beta}(-xy),\\
 p_{3\alpha+\beta}(y)p_\beta(x)
  ={}&p_\beta(x)p_{3\alpha+\beta}(y)p_{3\alpha+2\beta}(-xy).
\end{aligned}
\]
The corresponding relations in \(U'\) are
\[
\begin{aligned}
 p'_{\alpha'}(y')p'_{\beta'}(x')
  ={}&p'_{\beta'}(x')p'_{\alpha'}(y')
  p'_{\alpha'+\beta'}(-x'y')p'_{\alpha'+2\beta'}(-x'^2y')\\
  &{}\cdot p'_{\alpha'+3\beta'}(-x'^3y')
  p'_{2\alpha'+3\beta'}(-x'^3y'^2),\\
 p'_{\alpha'+\beta'}(y')p'_{\beta'}(x')
  ={}&p'_{\beta'}(x')p'_{\alpha'+\beta'}(y')
  p'_{\alpha'+2\beta'}(x'y'),\\
 p'_{\alpha'+3\beta'}(y')p'_{\alpha'}(x')
  ={}&p'_{\alpha'}(x')p'_{\alpha'+3\beta'}(y')
  p'_{2\alpha'+3\beta'}(x'y').
\end{aligned}
\]
Define \(\phi:U\to U'\) by
\[
\begin{aligned}
 &\phi\!\left(
 p_\alpha(x)p_\beta(y)p_{\alpha+\beta}(t)p_{2\alpha+\beta}(u)
 p_{3\alpha+\beta}(v)p_{3\alpha+2\beta}(w)\right)\\
 &\quad=
 p'_{\alpha'}(x^{3q})p'_{\beta'}(y^q)
 p'_{\alpha'+3\beta'}(t^{3q})p'_{2\alpha'+3\beta'}(u^{3q})
 p'_{\alpha'+\beta'}(v^q)p'_{\alpha'+2\beta'}(w^q).
\end{aligned}
\]
Substituting \(x^{3q},y^q\) in the three displayed relations for
\(U'\) yields exactly the images of the three relations for \(U\);
thus \(\phi\) is a homomorphism.

\SGARefTarget{sga3:XXIII:4.1.9}
\SGAParagraph{4.1.9. Semisimple rank greater than \(2\)}
For \(\alpha\in\Delta\), put \(\alpha'=d(\alpha)\). For every
\((\alpha,\beta)\), consider the pinned rank-at-most-two groups
\[
  Z_{\alpha\beta},\qquad Z'_{\alpha'\beta'}.
\]
The homomorphism \(M'\to M\) underlying \(h\) induces a
\(p\)-morphism
\[
  h_{\alpha\beta}:R(Z'_{\alpha'\beta'})
    \longrightarrow R(Z_{\alpha\beta}).
\]
By the preceding cases, there is a morphism of pinned groups
\[
  f_{\alpha\beta}:Z_{\alpha\beta}
    \longrightarrow Z'_{\alpha'\beta'}
\]
with \(R(f_{\alpha\beta})=h_{\alpha\beta}\). The restrictions
\[
  f_{\alpha\beta}|_{Z_\alpha},
  \qquad f_{\alpha\alpha}:Z_\alpha\to Z'_{\alpha'}
\]
correspond to the same morphism of pinned root data, and therefore
coincide by uniqueness. \SGARefLink{sga3:XXIII:2.4}{Corollary 2.4} now glues the
\(f_{\alpha\beta}\) to a morphism \(f:G\to G'\). It is a morphism of
pinned groups and satisfies \(R(f)=h\), completing the proof.
